% ============================================================
%  Modul 8 — Implementasi Multi-threading
% ============================================================

\chapter{Implementasi Multi-threading}
\label{chap:modul-08}

% --- 1. Tujuan dan Output Praktikum ---
\section{Tujuan dan Output Praktikum}

\subsection{Tujuan Praktikum}
Tujuan yang diharapkan dapat dicapai oleh mahasiswa adalah:
\begin{enumerate}
	\item Memahami arsitektur perangkat lunak dalam manajemen \textit{thread} (termasuk abstraksi \textit{Thread Control Block}) dan algoritma penjadwalan (\textit{scheduling}) pada tingkat inti sistem operasi (\textit{kernel-level}).
	\item Mempraktikkan implementasi logika \textit{Context Switching} untuk memfasilitasi eksekusi \textit{multi-threading} melalui penyimpanan dan pemuatan status register CPU (memanfaatkan manipulasi bahasa tingkat rendah/Assembly).
\end{enumerate}

\subsection{Output Praktikum}
Hasil luaran yang diekspektasikan adalah kelengkapan instrumen berikut:
\begin{itemize}
	\item Implementasi fungsional dari modifikasi berkas sumber (khususnya penyesuaian pada logika \texttt{thread.c}, \texttt{switch.asm}, dan \texttt{scheduler.h}) yang terbukti mampu mengompilasi rekayasa antrean \textit{thread}.
	\item Tangkapan layar (\textit{screenshot}) keluaran terminal dari emulator QEMU yang memvalidasi bahwa beberapa \textit{thread} (contoh: Thread A, B, dan C) sukses berjalan saling bergantian secara sinkron mengambil alih siklus CPU.
\end{itemize}

% --- 2. Dasar Teori ---
\newpage
\section{Dasar Teori}

\subsection{Pengenalan Thread dan Multi-threading}
Dalam sistem komputasi, \textit{thread} adalah unit instruksi terkecil yang dapat dijadwalkan oleh sistem operasi \cite{Silberschatz2018}. Jika sebuah proses diibaratkan sebagai program yang sedang berjalan dan memiliki memori sendiri, maka \textit{thread} adalah entitas eksekusi operasional di dalam proses tersebut. Lingkungan OS \textit{bare-metal} merintis kerangka \textit{multi-threading} agar dapat mensimulasikan penanganan banyak fungsi bahasa C secara bersamaan (konkurensi) \cite{pintos_stanford}. Secara konseptual, rutinitas pembuatan dan pengelolaan \textit{thread} dikonstruksi guna mendistribusikan beban kerja secara seimbang sehingga mendongkrak efisiensi utilitas prosesor.

\subsection{Konteks Eksekusi dan Context Switching}
Konteks eksekusi merujuk pada seluruh status kumpulan register dari prosesor (seperti EAX, EBX, EBP, dan ESP pada perangkat arsitektur x86) yang dikenakan oleh sebuah \textit{thread} pada momentum waktu tertentu \cite{bryant2015computer}. Ketika sistem operasi hendak memberhentikan satu \textit{thread} lalu beralih ke giliran antrean eksekusi milik \textit{thread} yang lain, OS wajib menukarkan keadaannya lewat taktik perpindahan \textit{Context Switching}.

Karena kerumitan modifikasi letak tumpukan memori (\textit{stack}) amat bergantung pada wujud arsitektur silikon sistem biner, prosedur krusial ini umumnya direpresentasikan murni melalui arahan instruksi pada level \textit{Assembly}. Operasi tersebut mutlak ditugaskan untuk mengekstrak simpanan (\textit{Push}) dari *thread* lama seraya mengangkat muatan register (\textit{Pop}) profil milik antrean *thread* pengganti.

\subsection{Stack Awal \textit{Thread} Baru}
Sebuah \textit{thread} baru tidak muncul dalam keadaan "siap jalan" secara ajaib. Ketika \texttt{thread\_create()} dipanggil, kernel harus menyiapkan \textit{stack} awal sintetis yang isinya menyerupai hasil penyimpanan konteks normal. Dengan cara ini, saat \textit{scheduler} pertama kali memilih \textit{thread} tersebut, kode \textit{context switch} cukup melakukan \textit{restore register} dan melompat ke alamat yang sudah disusun sebelumnya.

Pada praktiknya, \textit{stack} awal ini biasanya memuat alamat fungsi masuk, parameter bantu, dan beberapa nilai register awal. Karena itulah modul ini mengarahkan mahasiswa menyusun beberapa \textit{frame} berlapis di memori. Tujuannya bukan sekadar memenuhi format struktur data, melainkan memastikan perpindahan pertama menuju \textit{thread} baru berlangsung melalui jalur yang sama seperti perpindahan konteks biasa.

\subsection{Siklus Hidup \textit{Thread} dan \textit{Ready Queue}}
Sepanjang hidupnya, sebuah \textit{thread} akan berpindah-pindah status. Setelah dibuat, \textit{thread} masuk ke status \textit{Ready}, yaitu siap dijadwalkan tetapi belum memegang CPU. Ketika dipilih oleh \textit{scheduler}, statusnya berubah menjadi \textit{Running}. Jika ia menyerahkan CPU secara sukarela melalui \texttt{yield()}, ia kembali ke antrean \textit{Ready}. Jika menunggu suatu kejadian atau sumber daya, statusnya menjadi \textit{Blocked}. Setelah pekerjaannya selesai, \textit{thread} akan keluar dari siklus ini.

\textit{Ready queue} adalah antrean \textit{thread} yang siap dieksekusi. TCB menyimpan identitas dan konteks tiap \textit{thread}, sedangkan \textit{ready queue} menyimpan urutan siapa yang berhak dipertimbangkan oleh \textit{scheduler}. Pemisahan konsep ini penting: TCB menjawab "siapa \textit{thread}-nya dan apa keadaannya", sedangkan antrean menjawab "siapa yang menunggu giliran CPU".

Pada modul ini, \textit{scheduler} yang dibangun masih bersifat sederhana. Ia belum memakai \textit{timer interrupt} untuk merebut CPU secara paksa. Artinya, perpindahan antar \textit{thread} sangat bergantung pada titik-titik sukarela seperti \texttt{yield()} atau blokade karena sinkronisasi.


\subsection{\textit{Cooperative Round-Robin} pada Modul Ini}
Algoritma penjadwalan yang dipakai sebagai fondasi modul ini adalah \textit{cooperative round-robin}. Istilah \textit{round-robin} berarti \textit{thread} yang siap dijalankan diputar bergiliran dari antrean \textit{ready}. Istilah \textit{cooperative} berarti perpindahan hanya terjadi ketika \textit{thread} yang sedang berjalan rela melepaskan CPU atau masuk ke kondisi \textit{blocked}; bukan karena dipotong paksa oleh \textit{timer interrupt}.

Kelebihan pendekatan ini adalah implementasinya lebih mudah dipahami untuk tahap awal, karena mahasiswa dapat memusatkan perhatian pada struktur antrean, penyimpanan konteks, dan alur \texttt{yield()}. Namun, model ini juga memiliki keterbatasan: jika satu \textit{thread} tidak pernah menyerahkan CPU, maka \textit{thread} lain dapat menunggu sangat lama. Keterbatasan inilah yang nantinya menjadi jembatan menuju pembahasan sinkronisasi dan penjadwalan yang lebih maju pada modul-modul berikutnya.

\subsection{Thread Control Block / TCB}
Agar sistem pengontrol sanggup menjaga identitas serta alokasi titik perhentian terakhir masing-masing \textit{thread}, direkayasa sebuah skema pelacak terstruktur dinamakan \textit{Thread Control Block} (TCB) \cite{Silberschatz2018}. Menggunakan struktur data mandiri (seperti halnya tipe rekaman mendasar pada bahasa C), blok TCB secara dedikatif membungkus parameter-parameter memori berikut:
\begin{itemize}
	\item \textbf{Thread ID}: Angka unit unik pengenal subyek yang tak boleh identik.
	\item \textbf{Status Thread}: Peta kondisi berjalannya instruksi pada suatu \textit{thread} (contoh: status sedang fasif siap diterjunkan/\textit{Ready}, status mandek/\textit{Blocked}, atau sedang berjalan/\textit{Running}).
	\item \textbf{Stack Pointer (ESP)}: Penunjuk referensi baris alamat terakhir yang terekam pada tumpukan (\textit{stack}) ketika interupsi perpindahan prosesnya berlangsung.
\end{itemize}

\subsection{Algoritma Penjadwalan (Scheduling)}
Siklus perancangan eksekusi CPU dikoordinasi secara tersentralisasi oleh modul OS benama \textit{Scheduler} \cite{Silberschatz2018}. Esensi utama dari algoritma komponen ini bertugas menyeleksi manakah prioritas \textit{thread} yang menduduki gerbang \textit{Ready} pada antrean sistem, agar mendapat wewenangnya dialihkan ke tingkatan fase \textit{Running}. Konsep paradigma fundamental penyeleksian dijabarkan melalui dua cara:
\begin{enumerate}
	\item \textbf{Cooperative Scheduling}: Sebuah \textit{thread} secara sadar membiarkan algoritma lain ganti mengeksekusi CPU dengan inisiatifnya menyerah (seperti metode inkuiri \texttt{yield()} atau tertahan menunggu sumber daya antarmuka selesai).
	\item \textbf{Preemptive Scheduling}: Sistem operasi menyetel batasan pencacah peranti keras layaknya \textit{Timer Interrupt}, berfungsi untuk merampas status dominasi algoritma tanpa ampun dari instruksi \textit{thread} yang sedang menduduki memori apakala batasan waktunya telah dirasa mencukupi.
\end{enumerate}

Tabel \ref{tab:m08-dasar-teori} memberikan intisari terminologi vital terkait landasan teori \textit{multithreading}.

\begin{table}[!ht]
	\caption{Tabel Ringkasan Dasar Teori Multi-threading}
	\label{tab:m08-dasar-teori}
	\centering
	\begin{tabulary}{\textwidth}{|L|L|}
		\hline
		\textbf{Konsep} & \textbf{Penjelasan Singkat} \\ \hline
		Thread & Unit eksekusi fungsional terkecil di dalam OS. \\ \hline
		Context Switch & Tindakan rekam-simpan silang arsitektur register satu \textit{thread} dengan \textit{thread} cadangan lainnya. \\ \hline
		TCB & Struktur blok tipe data spesifik pengingat status profil \textit{thread}. \\ \hline
		Scheduler & Organ algoritma OS yang mendikte hak pembagian waktu kendali CPU atas \textit{thread}. \\ \hline
	\end{tabulary}
\end{table}

% --- 3. Alat dan Bahan ---
\newpage
\section{Alat dan Bahan}
\label{sec:m08-alat-bahan}

\subsection{Perangkat Keras}
\begin{itemize}
	\item PC atau Laptop dengan arsitektur prosesor basis \textit{host} x86/x64 yang mendukung kelanjutan kompilasi emulator.
	\item Memori minimal 4GB RAM ruang eksekusi.
\end{itemize}

\subsection{Perangkat Lunak (Bahan Praktikum)}
\begin{itemize}
	\item OS \textit{Host} utama (Windows dengan WSL/GitBash, Linux, atau macOS).
	\item Text Editor (misalnya Visual Studio Code) untuk menyunting berkas kode bahasa C dan \textit{Assembly}.
	\item Instalasi Docker Engine (\texttt{docker.io}) untuk menyalakan lingkungan isolasi kompilasi.
	\item Emulator piranti keras x86: \texttt{qemu-system-x86} (atau instalasi QEMU setara) yang akan mengeksekusi binari hasil rakitan OS.
	\item Direktori \textit{Starter Code} OS 32-bit yang sudah disiapkan sebelumnya (berisi direktori \texttt{src} dengan kerangka lengkap \texttt{kernel.c}, \texttt{thread.h}, \texttt{thread.c}, \texttt{switch.asm}, \texttt{scheduler.h}, dsb).
\end{itemize}

% --- 4. Langkah-Langkah Praktikum ---
\newpage
\section{Langkah-Langkah Praktikum}

\subsection{Sesi 1: Persiapan Lingkungan Kerja}
\label{sec:m08-sesi-1}
Langkah pertama adalah menyiapkan direktori \textit{Starter Code} dan masuk ke dalam kontainer Docker.

\begin{enumerate}
	\item Buka terminal dan pastikan Anda berada di dalam folder \texttt{starter-code-os32-main}.
	\item Masuk ke lingkungan Docker dengan menjalankan perintah:
	      \begin{lstlisting}[language=bash, caption={Masuk ke lingkungan kontainer}]
./run-docker.sh
    \end{lstlisting}
	\item Buka berkas \texttt{src/thread.c} pada teks editor. Berkas ini adalah file utama yang akan dimodifikasi.
\end{enumerate}

\subsection{Sesi 2: Implementasi Fungsi \texttt{thread\_create()}}
\label{sec:m08-sesi-2}
Fungsi \texttt{thread\_create()} bertugas menyiapkan \textit{Thread Control Block} (TCB) dan \textit{stack} awal untuk \textit{thread} baru. Buka \texttt{thread.c}, cari fungsi pengisian \texttt{thread\_create}, dan ikuti langkah berikut:

\begin{enumerate}
	\item \textbf{Langkah 1 (Alokasi TCB):} Gunakan \texttt{allocate\_thread\_page()} untuk alokasi memori \textit{thread}. Jika berhasil, panggil \texttt{init\_thread(t, name, priority)} dan berikan ID dengan \texttt{t->tid = allocate\_tid()}.
	      \begin{figure}[H]
		      \centering
		      \includegraphics[width=0.8\textwidth]{Figure/08/1.png}
		      \caption{Implementasi Alokasi TCB pada fungsi thread\_create()}
	      \end{figure}

	\item \textbf{Langkah 2 (Frame 1: \texttt{kernel\_thread\_frame}):} Buat variabel \textit{frame} ini menggunakan \texttt{alloc\_frame()}. Isi \texttt{eip} dengan \texttt{NULL}, lalu isi \texttt{function} dan \texttt{aux} sesuai dengan argumen parameter yang dikirim.
	      \begin{figure}[H]
		      \centering
		      \includegraphics[width=0.8\textwidth]{Figure/08/2.png}
		      \caption{Implementasi Frame 1 pada fungsi thread\_create()}
	      \end{figure}

	\item \textbf{Langkah 3 (Frame 2: \texttt{switch\_entry\_frame}):} Buat variabel \textit{frame} ini dengan alokasi di bawah Frame 1. Isi \texttt{eip} dengan memanggil alamat fungsi \texttt{kernel\_thread}.
	      \begin{figure}[H]
		      \centering
		      \includegraphics[width=0.8\textwidth]{Figure/08/3.png}
		      \caption{Implementasi Frame 2 pada fungsi thread\_create()}
	      \end{figure}

	\item \textbf{Langkah 4 (Frame 3: \texttt{switch\_threads\_frame}):} Buat variabel \textit{frame} terbawah ini. Isi variabel \texttt{eip} dengan \texttt{switch\_entry}, lalu set nilai \texttt{ebp}, \texttt{ebx}, \texttt{esi}, dan \texttt{edi} menjadi \texttt{0}.
	      \begin{figure}[H]
		      \centering
		      \includegraphics[width=0.8\textwidth]{Figure/08/4.png}
		      \caption{Implementasi Frame 3 pada fungsi thread\_create()}
	      \end{figure}

	\item \textbf{Langkah 5 (Penyelesaian):} Panggil \texttt{thread\_unblock(t)} agar dimasukkan ke dalam antrean \textit{ready}, kemudian lakukan \textit{return} id \textit{thread} tersebut (\texttt{t->tid}).
	      \begin{figure}[H]
		      \centering
		      \includegraphics[width=0.8\textwidth]{Figure/08/5.png}
		      \caption{Langkah Penyelesaian pada fungsi thread\_create()}
	      \end{figure}

	\item Pastikan kerangka utuh fungsi Anda menyerupai format referensi ini:
	      \begin{figure}[H]
		      \centering
		      \includegraphics[width=0.8\textwidth]{Figure/08/6.png}
		      \caption{Fungsi utuh \texttt{thread\_create()} secara keseluruhan}
	      \end{figure}
\end{enumerate}

\subsection{Sesi 3: Implementasi Cooperative \textit{Round-Robin}}
\label{sec:m08-sesi-3}
Bagian ini mengatur agar \textit{thread} dapat saling bergantian menggunakan wewenang CPU. Masih di dalam file \texttt{thread.c}:

\begin{enumerate}
	\item \textbf{Modifikasi \texttt{thread\_yield()}:} Ambil \textit{thread} yang saat ini berjalan menggunakan \texttt{thread\_current()}. Ubah statusnya menjadi \texttt{THREAD\_READY}. Masukkan \textit{thread} ini ke dalam antrean \texttt{ready\_queue} pada indeks \texttt{ready\_tail}. Geser urutan \texttt{ready\_tail}, tambahkan jumlah \texttt{ready\_count}, dan akhiri dengan memanggil fungsi \texttt{schedule()}.
	      \begin{figure}[H]
		      \centering
		      \includegraphics[width=0.8\textwidth]{Figure/08/7.png}
		      \caption{Implementasi fungsi \texttt{thread\_yield()}}
	      \end{figure}

	\item \textbf{Modifikasi \texttt{next\_thread\_to\_run()}:} Fungsi ini berjalan secara metode FIFO (\textit{First-In-First-Out}). Jika antrean kosong (\texttt{ready\_count} bernilai 0), kembalikan \texttt{initial\_thread}. Jika isi antrean tersedia, ambil \textit{thread} dari urutan depan (\texttt{ready\_queue[ready\_head]}), geser letak indeks depan \texttt{ready\_head}, kurangi nilai hitung \texttt{ready\_count}, lalu \textit{return thread} yang diambil tersebut.
	      \begin{figure}[H]
		      \centering
		      \includegraphics[width=0.8\textwidth]{Figure/08/8.png}
		      \caption{Implementasi fungsi \texttt{next\_thread\_to\_run()}}
	      \end{figure}
\end{enumerate}

\subsection{Sesi 4: Kompilasi dan Pengujian QEMU}
\label{sec:m08-sesi-4}
Setelah seluruh kode di \texttt{thread.c} selesai, jalankan pengujian sistem untuk membuktikan bahwa OS sudah sanggup menjalankan perpindahan waktu \textit{thread}.

\begin{enumerate}
	\item Buka kembali terminal yang berada di dalam lingkungan Docker, lalu jalankan perintah kompilasi:
	      \begin{lstlisting}[language=bash, caption={Perintah build pada terminal Docker}]
make
    \end{lstlisting}
	\item Setelah proses kompilasi selesai (\textit{error} tidak muncul), jalankan QEMU emulator secara langsung menggunakan perintah di bawah ini:
	      \begin{lstlisting}[language=bash, caption={Menjalankan OS pada terminal Docker}]
make run
    \end{lstlisting}
	\item Perhatikan layar simulator QEMU. Anda harus masuk dalam putaran sistem \textit{infinite loop} yang memperlihatkan ketiga \textit{thread} yang dibuat mencetak teks (Thread A, Thread B, dan Thread C) bergantian per waktu putaran:
	      \begin{lstlisting}[caption={Ekspektasi output sistem penjadwalan Thread di latar QEMU}]
===================================
  IF-OS  |  Praktikum Sistem Operasi
===================================

Inisialisasi sistem thread...
Membuat thread A, B, C...
Memulai scheduling...

[Thread A] tick 0
[Thread B] tick 0
[Thread C] tick 0
[Thread A] tick 1
[Thread B] tick 1
...
    \end{lstlisting}
\end{enumerate}

% ----------------------------------------------------------------
\section{Latihan}

% ----------------------------------------------------------------
\subsection{Tujuan}

Jawablah pertanyaan-pertanyaan berikut berdasarkan hasil praktikum Anda
sesuai urutan sesi pada Modul~8.
Sertakan bukti berupa \textit{screenshot} atau penggalan keluaran
terminal yang relevan pada setiap jawaban.

% ----------------------------------------------------------------
\subsection{Sesi 1: Persiapan Lingkungan Kerja}

\subsubsection{Masuk ke Container Docker}

Jalankan perintah berikut di terminal, kemudian masuk ke direktori
\textit{workspace}:

\begin{lstlisting}[language=bash, caption={Menjalankan container dan masuk workspace}]
docker run --rm -it \
    -v "C:\path\ke\starter-code:/workspace" \
    sinanonym/ifos-env
cd /workspace
\end{lstlisting}

\begin{enumerate}

	\item \textbf{Compile starter code sebelum implementasi.}\\
	      Jalankan perintah berikut di dalam container:

	      \begin{lstlisting}[language=bash, caption={Perintah kompilasi awal}]
make clean && make
    \end{lstlisting}

	      Tuliskan seluruh \textit{warning} yang muncul pada saat kompilasi dan
	      jelaskan mengapa \textit{warning} tersebut muncul pada kode
	      \textit{stub}.

	\item \textbf{Jalankan QEMU sebelum implementasi.}\\
	      Jalankan perintah berikut lalu sertakan \textit{screenshot} layar QEMU.
	      Jelaskan mengapa sistem berhenti di sana dan belum menampilkan
	      output Thread~A/B/C.

	      \begin{lstlisting}[language=bash, caption={Menjalankan emulator QEMU}]
make run
    \end{lstlisting}


\end{enumerate}

% ----------------------------------------------------------------
\subsection{Sesi 2: Implementasi \texttt{thread\_create()}}

\subsubsection{Thread Control Block dan Layout Memori}

\begin{enumerate}

	\item \textbf{Field-field \texttt{struct thread}.}\\
	      Perhatikan definisi \texttt{struct thread} di \texttt{src/thread.h}.
	      Isilah tabel berikut berdasarkan kode sumber.

	      \begin{table}[!ht]
		      \caption{Field-field pada \texttt{struct thread}}
		      \label{tab:h04-struct-thread}
		      \centering
		      \begin{tabulary}{\linewidth}{|L|L|L|}
			      \hline
			      \textbf{Field} & \textbf{Tipe Data} & \textbf{Fungsi} \\
			      \hline
			      \texttt{tid}      & \texttt{tid\_t}              & \textit{[Isi]} \\
			      \hline
			      \texttt{status}   & \texttt{enum thread\_status} & \textit{[Isi]} \\
			      \hline
			      \texttt{name}     & \texttt{char[16]}            & \textit{[Isi]} \\
			      \hline
			      \texttt{stack}    & \texttt{uint8\_t *}          & \textit{[Isi]} \\
			      \hline
			      \texttt{priority} & \texttt{int}                 & \textit{[Isi]} \\
			      \hline
			      \texttt{magic}    & \texttt{unsigned}            & \textit{[Isi]} \\
			      \hline
		      \end{tabulary}
	      \end{table}

	\item \textbf{Layout halaman 4~KB setiap thread.}\\
	      Setiap thread menempati satu halaman (\textit{page}) memori berukuran
	      4~KB. Isilah kolom \textit{Isi} pada tabel layout di bawah ini
	      untuk menggambarkan posisi \texttt{struct thread}, \textit{stack},
	      dan ketiga frame yang disiapkan oleh \texttt{thread\_create()}.

	      \begin{table}[!ht]
		      \caption{Layout halaman 4~KB milik sebuah thread}
		      \label{tab:h04-page-layout}
		      \centering
		      \begin{tabulary}{\linewidth}{|C|L|}
			      \hline
			      \textbf{Offset (relatif)} & \textbf{Isi} \\
			      \hline
			      \texttt{+4096} & \textit{[Isi: puncak halaman / ESP awal]} \\
			      \hline
			      \texttt{...}   & \texttt{switch\_threads\_frame} \\
			      \hline
			      \texttt{...}   & \texttt{switch\_entry\_frame} \\
			      \hline
			      \texttt{...}   & \texttt{kernel\_thread\_frame} \\
			      \hline
			      \texttt{...}   & \textit{[Isi: area stack kosong, tumbuh ke bawah ($\downarrow$)]} \\
			      \hline
			      \texttt{+0}    & \texttt{struct thread} (TCB) \\
			      \hline
		      \end{tabulary}
	      \end{table}

	\item \textbf{Analisis fungsi \texttt{running\_thread()}.}\\
	      Perhatikan potongan kode berikut:

	      \begin{lstlisting}[language=c, caption={Implementasi \texttt{running\_thread()} di \texttt{src/thread.c}}]
struct thread *running_thread(void) {
    uint32_t esp;
    __asm__ volatile ("mov %%esp, %0" : "=g" (esp));
    return (struct thread *)(void *)(esp & ~(uint32_t)(PGSIZE - 1));
}
    \end{lstlisting}

	      \begin{enumerate}
		      \item[(a)] Jelaskan apa yang dilakukan operasi bitwise
		            \texttt{esp \& \textasciitilde(PGSIZE - 1)} dan
		            mengapa hasilnya selalu merupakan alamat awal
		            \textit{page}.

		      \item[(b)] Mengapa teknik ini hanya benar jika
		            \texttt{struct thread} selalu berada di offset~0
		            dari halamannya? Apa dampaknya jika
		            \texttt{initial\_thread\_page} di \texttt{boot.asm}
		            tidak di-\textit{align} ke 4096~byte?

		      \item[(c)] Apa fungsi \texttt{THREAD\_MAGIC} dan bagaimana
		            \texttt{thread\_current()} menggunakannya untuk
		            mendeteksi \textit{stack overflow}?
	      \end{enumerate}

\end{enumerate}

\subsubsection{Implementasi Fungsi \texttt{thread\_create()}}

\begin{enumerate}

	\item \textbf{Tiga stack frame pada \texttt{thread\_create()}.}\\
	      Isilah tabel berikut berdasarkan komentar petunjuk di
	      \texttt{src/thread.c}.

	      \begin{table}[!ht]
		      \caption{Tiga stack frame yang disiapkan \texttt{thread\_create()}}
		      \label{tab:h04-stack-frames}
		      \centering
		      \begin{tabulary}{\linewidth}{|C|L|L|L|}
			      \hline
			      \textbf{Urutan push} & \textbf{Nama frame} &
			      \textbf{Field penting} & \textbf{Nilai yang diisi} \\
			      \hline
			      1, paling atas & \texttt{kernel\_\allowbreak thread\_\allowbreak frame} &
			      \texttt{eip}, \texttt{function}, \texttt{aux} &
			      \textit{[Isi]} \\
			      \hline
			      2, tengah & \texttt{switch\_\allowbreak entry\_\allowbreak frame} &
			      \texttt{eip} &
			      \textit{[Isi]} \\
			      \hline
			      3, paling bawah & \texttt{switch\_\allowbreak threads\_\allowbreak frame} &
			      \texttt{eip}, register \textit{callee-saved} &
			      \textit{[Isi]} \\
			      \hline
		      \end{tabulary}
	      \end{table}

	\item \textbf{Alur eksekusi pertama thread baru.}\\
	      Saat thread baru \textbf{pertama kali} dijadwalkan melalui
	      \texttt{switch\_threads()}, jelaskan secara berurutan apa yang
	      terjadi mulai dari \texttt{switch\_entry} hingga fungsi thread
	      (\texttt{func\_a}, \texttt{func\_b}, atau \texttt{func\_c})
	      benar-benar mulai berjalan.

	\item \textbf{Tulis implementasi \texttt{thread\_create()}.}\\
	      Salin kode yang telah Anda tulis di \texttt{src/thread.c}:

	      \begin{lstlisting}[language=c, caption={Implementasi \texttt{thread\_create()} oleh mahasiswa}]
tid_t thread_create(const char *name, int priority,
                    thread_func *function, void *aux) {

    /* Tulis implementasi Anda di sini */

}
    \end{lstlisting}

	\item \textbf{Penanganan kegagalan alokasi.}\\
	      Kapan \texttt{allocate\_thread\_page()} dapat mengembalikan
	      \texttt{NULL}, dan bagaimana \texttt{thread\_create()} harus
	      menanganinya?

\end{enumerate}

% ----------------------------------------------------------------
\subsection{Sesi 3: Implementasi \textit{Cooperative Round-Robin}}

\subsubsection{Konsep Status Thread}

\begin{enumerate}

	\item \textbf{Perbedaan \texttt{THREAD\_READY} dan
		      \texttt{THREAD\_BLOCKED}.}\\
	      Lengkapi tabel perbandingan berikut.

	      \begin{table}[!ht]
		      \caption{Perbandingan status \texttt{THREAD\_READY} dan \texttt{THREAD\_BLOCKED}}
		      \label{tab:h04-thread-status}
		      \small
		      \begin{tabularx}{\linewidth}{|X|X|X|}
			      \hline
			      \textbf{Aspek} &
			      \makecell[l]{\textbf{THREAD\_}\\\textbf{READY}} &
			      \makecell[l]{\textbf{THREAD\_}\\\textbf{BLOCKED}} \\
			      \hline
			      Siapa yang menetapkan? & \textit{[Isi]} & \textit{[Isi]} \\
			      \hline
			      Masuk ke \textit{ready queue}? & \textit{[Isi]} & \textit{[Isi]} \\
			      \hline
			      Cara kembali ke \texttt{RUNNING} & \textit{[Isi]} & \textit{[Isi]} \\
			      \hline
			      Contoh penggunaan & \textit{[Isi]} & \textit{[Isi]} \\
			      \hline
		      \end{tabularx}
	      \end{table}

\end{enumerate}

\subsubsection{Implementasi Fungsi \textit{Scheduling}}

\begin{enumerate}

	\item \textbf{Tulis implementasi \texttt{thread\_yield()}.}\\
	      Salin kode yang telah Anda tulis di \texttt{src/thread.c}:

	      \begin{lstlisting}[language=c, caption={Implementasi \texttt{thread\_yield()} oleh mahasiswa}]
void thread_yield(void) {

    /* Tulis implementasi Anda di sini */

}
    \end{lstlisting}

	\item \textbf{Tulis implementasi \texttt{next\_thread\_to\_run()}
		      (Round-Robin FIFO).}\\
	      Salin kode yang telah Anda tulis di \texttt{src/thread.c}:

	      \begin{lstlisting}[language=c, caption={Implementasi \texttt{next\_thread\_to\_run()} oleh mahasiswa}]
struct thread *next_thread_to_run(void) {

    /* Tulis implementasi Anda di sini */

}
    \end{lstlisting}

	\item \textbf{Mengapa mengembalikan \texttt{initial\_thread},
		      bukan \texttt{NULL}?}\\
	      Jelaskan apa yang akan terjadi pada fungsi \texttt{schedule()} jika
	      \texttt{next\_thread\_to\_run()} mengembalikan \texttt{NULL} saat
	      antrian kosong.

	\item \textbf{\texttt{thread\_block()} vs.\ \texttt{thread\_yield()}
		      di akhir \texttt{kernel\_main}.}\\
	      Jelaskan perbedaan perilaku kedua fungsi tersebut setelah Thread~A,
	      B, dan C berhasil dibuat, dan mengapa \texttt{thread\_block()}
	      adalah pilihan yang benar.

\end{enumerate}

% ----------------------------------------------------------------
\subsection{Sesi 4: Kompilasi dan Pengujian QEMU}

\begin{enumerate}

	\item \textbf{Compile ulang setelah implementasi selesai.}

	      \begin{lstlisting}[language=bash, caption={Perintah build dan run setelah implementasi}]
make clean && make && make run
    \end{lstlisting}

	      Sertakan \textit{screenshot} layar QEMU yang menampilkan ketiga
	      thread berjalan bergantian.


	\item \textbf{Verifikasi urutan penjadwalan.}\\
	      Ekspektasi keluaran sistem yang benar adalah sebagai berikut:

	      \begin{lstlisting}[caption={Ekspektasi output Round-Robin di layar QEMU}]
===================================
  IF-OS  |  Praktikum Sistem Operasi
===================================

Inisialisasi sistem thread...
Membuat thread A, B, C...
Memulai scheduling...

[Thread A] tick 0
[Thread B] tick 0
[Thread C] tick 0
[Thread A] tick 1
[Thread B] tick 1
[Thread C] tick 1
...
    \end{lstlisting}

	      Apakah urutan Thread~A $\to$ Thread~B $\to$ Thread~C $\to$
	      Thread~A $\to \ldots$ terjaga secara konsisten pada hasil Anda?
	      Jelaskan mengapa urutan ini dapat dijamin oleh implementasi FIFO
	      pada \textit{ready queue}.

	\item \textbf{Analisis \textit{warning} compiler sebelum dan setelah
		      implementasi.}\\
	      Apakah \textit{warning} \texttt{unused-function} dan
	      \texttt{unused-variable} yang muncul sebelum implementasi masih
	      ada setelah Anda mengimplementasikan ketiga fungsi? Jelaskan
	      mengapa \textit{warning} tersebut hilang.

\end{enumerate}

% ----------------------------------------------------------------
\subsection{Eksplorasi Mandiri}

\begin{enumerate}

	\item \textbf{Menambah Thread keempat.}\\
	      Tambahkan Thread~D dengan fungsi dan warna VGA berbeda di
	      \texttt{src/kernel.c}, kemudian compile ulang dan jalankan.

	      \begin{enumerate}
		      \item[(a)] Sertakan \textit{screenshot} layar QEMU yang
		            menampilkan empat thread berjalan bergantian.


		      \item[(b)] Berapa nilai \texttt{MAX\_THREADS} saat ini?
		            Apa yang terjadi jika Anda membuat thread melebihi
		            batas tersebut? Tunjukkan dengan percobaan dan
		            sertakan bukti.
	      \end{enumerate}

	\item \textbf{Simulasi \textit{finite thread}.}\\
	      Ubah \texttt{func\_a} di \texttt{src/kernel.c} agar hanya
	      mencetak \textbf{5 kali} lalu berhenti (ganti \texttt{while(1)}
	      dengan loop terbatas). Amati apakah Thread~B dan Thread~C tetap
	      berjalan setelah Thread~A selesai, lalu jelaskan peran
	      \texttt{thread\_exit()} dalam skenario ini.

	\item \textbf{Refleksi: \textit{cooperative} vs.\
		      \textit{preemptive scheduling}.}\\
	      Sistem IF-OS saat ini menggunakan \textit{cooperative scheduling}
	      di mana thread hanya berganti saat secara sadar memanggil
	      \texttt{thread\_yield()}.

	      \begin{enumerate}
		      \item[(a)] Apa kelemahan utama \textit{cooperative scheduling}
		            dibandingkan \textit{preemptive scheduling}?
		            Berikan contoh skenario nyata di mana kelemahan
		            ini berbahaya.

		      \item[(b)] Komponen perangkat keras apa yang perlu ditambahkan
		            agar kernel dapat memaksa pergantian thread secara
		            otomatis tanpa bergantung pada kepatuhan thread
		            itu sendiri?
	      \end{enumerate}

\end{enumerate}
